\section{Worked Examples}
\label{sec:examples}

We illustrate $\Mtype\langle T \rangle$ on three fabricated scenarios. The
first is anchored in real code shipped with the
crate~\cite{hansen_deepcausality}; the second and third are design
sketches at the same abstraction level as the type theory
in~\Cref{sec:type}, included to show how the same pattern carries to
other sparse-data domains. We mark the status of each example explicitly.

\subsection{Aspirin headache trial}
\label{sec:examples:trial}

\paragraph{Status.} Fully worked example. Anchored in
\texttt{examples/clinical\_trial/clinical\_trial.rs} in the crate.

\paragraph{Setup.} A small fabricated trial compares aspirin against a
control placebo on five patients, each contributing a pain-reduction
report. The reports are deliberately heterogeneous: some patients report
every day with a tight value, some report intermittently with a noisy
value, and at least one patient produces no data at all. The four
constructors of $\Mtype\langle f64 \rangle$ map onto the situations
one-to-one. Patient A (control) is a known fact and uses
$\textsc{fromValue}(0.5)$. Patient B (strong responder) reports every day
with noisy value and uses $\textsc{fromUncertain}$ over a $\mathcal{N}(6,2)$
distribution. Patient C (weak intermittent responder) has presence
probability $0.3$ and uses $\textsc{fromBernoulli}(0.3,\,\mathcal{N}(1.0,
0.5))$. Patient D (moderate good-reporter) uses
$\textsc{fromBernoulli}(0.8,\,\mathcal{N}(4.0, 2.5))$. Patient E (control,
missing data) uses $\textsc{alwaysNone}$.

\paragraph{Behaviour under the algebra.} Adding two patient values is the
canonical use of the AND rule on presence:
$$
\Mtype\text{-sum}(A, B) = \bigl(\,x_A \wedge x_B,\; v_A + v_B\,\bigr).
$$
For patients A and B, both presence axes are $\delta_{\text{true}}$, so
the result is always present. For patients A and C, the result is present
with probability $0.3$, since A is certainly present and C carries the
Bernoulli factor. Sampling $\Mtype\text{-sum}(A, C)$ returns $\bot$ roughly
$70\%$ of the time, and a real value sampled from the sum distribution the
remaining $30\%$. The AND-on-presence rule prevents the sum from
synthesising a numeric value for a patient who did not report.

\paragraph{Decision via the gating operator.} The example demonstrates the
strict-versus-lax operating points of $\textsc{lift}$ on the same data.
Lifting patient D with $(p_\tau, 1-\alpha, \varepsilon, N) = (0.7, 0.95,
0.05, 1000)$ succeeds because the true presence probability $0.8$ is well
above the threshold; the result is a plain $\Utype\langle f64 \rangle$
whose mean can be inspected with $1000$ samples. Lifting patient C with
threshold $0.8$ fails with $\mathrm{PresenceError}$ because the true
presence probability $0.3$ is far below the threshold; SPRT accepts $H_0$
quickly. Lifting patient E with any threshold strictly greater than zero
fails for the same reason. The example then aggregates only the
successfully-lifted aspirin-group reductions, computes their average via
$\Utype$ division, and compares against the control via the
\texttt{greater\_than} comparator on $\Utype\langle f64 \rangle$. The
decision is made with $\texttt{probability\_exceeds}(0.9, 0.95, 0.05,
1000)$ on the result.

\paragraph{What the example shows.} Three things. First, the four
constructors are sufficient to encode a realistic mix of patient
reporting behaviours. Second, $\textsc{lift}$ functions as a programmable
inclusion criterion: the threshold and confidence parameters are the
trial designer's knobs for how strict to be about which patients
contribute to the aggregate. Third, the AND rule on presence is the
right default for arithmetic on intermittent data; the sum of a complete
report and a missing report is missing, not zero and not the
complete-report value.

\subsection{Intermittent sensor}
\label{sec:examples:sensor}

\paragraph{Status.} Design sketch. The crate's
\texttt{examples/sensor/sensor\_processing.rs} demonstrates
$\Utype\langle T \rangle$ on noisy sensors but does not exercise
$\Mtype\langle T \rangle$. The pattern below is the natural extension.

\paragraph{Setup.} An embedded sensor produces a noisy reading every
$\Delta t$ seconds; when the sensor is healthy a reading $r \sim
\mathcal{N}(\mu, \sigma^2)$ arrives, but on each tick the sensor has
probability $1 - p$ of dropping the reading. The application computes a
moving average of the last $k$ readings and triggers a downstream alert
if the average crosses a threshold.

\paragraph{Three baselines.} Three encodings of the same sensor are
useful to compare. The first treats every tick as a measurement and
substitutes a fallback value (often zero or a recent mean) on drop; this
is the imputation-laundering baseline that injects fabricated numbers
into the moving average. The second wraps each reading in
\texttt{Option<Uncertain<f64>>} and only includes \texttt{Some} readings;
this is the hard-gated-drop baseline that collapses presence into a
deterministic boolean. The third uses $\Mtype\langle f64 \rangle$ with
each reading constructed by $\textsc{fromBernoulli}(p,\,\mathcal{N}(\mu,
\sigma^2))$.

\paragraph{Behaviour under the algebra.} The moving average is a sum of
$k$ readings divided by $k$. Under the $\Mtype$ encoding, the AND rule on
presence makes the moving average present only if every tick in the
window was present; otherwise the moving average is $\bot$. The alert
condition can then be expressed as a presence-gated comparison: lift the
moving average to $\Utype\langle f64 \rangle$ at a threshold of $p^k -
\varepsilon$, and compare to the alert level using
\texttt{probability\_exceeds}. The window-presence threshold $p^k$
follows from independence of the per-tick Bernoulli draws and gives the
designer a single statistical parameter to control sensitivity in the
sparse regime.

\paragraph{What this sketch shows.} The presence axis combined with the
gating operator gives a natural place to encode the dropout-rate
assumption explicitly. The imputation baseline buries it in the choice of
fallback value; the hard-gated baseline buries it in the choice of
threshold for $k$. With $\Mtype$, the assumption appears in the
$\textsc{fromBernoulli}$ parameter and is propagated through the moving
average automatically.

\subsection{GPS dropouts}
\label{sec:examples:gps}

\paragraph{Status.} Design sketch. The crate's
\texttt{examples/gps/gps\_navigation.rs} demonstrates $\Utype\langle T
\rangle$ on noisy GPS readings but does not exercise $\Mtype\langle T
\rangle$. The pattern below is the natural extension.

\paragraph{Setup.} A vehicle's GPS receiver produces a fix every second.
When the receiver has line-of-sight the fix is noisy but valid; in
tunnels and urban canyons the receiver loses the fix entirely. A
navigation application computes speed from consecutive fixes by the
finite difference $v = \Delta d / \Delta t$.

\paragraph{The bug class to prevent.} Computing speed from a pair of
fixes one of which is missing is the central failure mode the example
addresses. Three encodings give three different bugs. Substituting the
last known fix for the missing one synthesises a zero speed and reports
``vehicle stopped'' even though the vehicle is moving inside a tunnel.
Dropping the pair entirely produces silent gaps in the speed trace and
leaves the user with no signal that anything was missing. Reporting a
``best effort'' speed by extrapolation invents both a value and a
confidence that the data does not support.

\paragraph{The $\Mtype$ encoding.} Each fix is a $\Mtype\langle f64
\rangle$ value (one for latitude, one for longitude, or a struct of two
$\Mtype$ values), with presence governed by a Bernoulli factor learned
from the historical dropout rate. The finite difference $\Delta d$ is a
subtraction between two $\Mtype$ fixes; under the AND rule the result is
present only when both fixes were present. Dividing by $\Delta t$
preserves the presence axis, so $v = \Delta d / \Delta t$ is a
$\Mtype\langle f64 \rangle$ that is present exactly when both fixes were
present. The application can then decide what to do with absent speed
samples: skip the update, raise an alert, fall back to a separate
inertial estimate, or expose the presence probability through
$\textsc{isSome}$ and let downstream code decide.

\paragraph{What this sketch shows.} The same pattern carries: the
presence axis records the structural fact that an output depends on the
existence of inputs, the AND rule propagates that dependence through
arithmetic without further code, and the gating operator gives the
application a principled place to make decisions about partial data.
